% offset-mechanism-equations.tex
%
% v2: reduced to the two corrections and nothing else. The full difference (Delta M,
%     the absorber column, K=0 and the two closed forms) moved to
%     offset-mechanism-derivation.tex, which this file does not depend on.
% v1: contained both; the mixture obscured the simple statement.
%
% THIS IS A FRAGMENT, not a standalone document. It uses \mhr from
% jan-augmentation-writeup.tex and is meant to be pulled in with
%     % offset-mechanism-equations.tex
%
% v2: reduced to the two corrections and nothing else. The full difference (Delta M,
%     the absorber column, K=0 and the two closed forms) moved to
%     offset-mechanism-derivation.tex, which this file does not depend on.
% v1: contained both; the mixture obscured the simple statement.
%
% THIS IS A FRAGMENT, not a standalone document. It uses \mhr from
% jan-augmentation-writeup.tex and is meant to be pulled in with
%     % offset-mechanism-equations.tex
%
% v2: reduced to the two corrections and nothing else. The full difference (Delta M,
%     the absorber column, K=0 and the two closed forms) moved to
%     offset-mechanism-derivation.tex, which this file does not depend on.
% v1: contained both; the mixture obscured the simple statement.
%
% THIS IS A FRAGMENT, not a standalone document. It uses \mhr from
% jan-augmentation-writeup.tex and is meant to be pulled in with
%     % offset-mechanism-equations.tex
%
% v2: reduced to the two corrections and nothing else. The full difference (Delta M,
%     the absorber column, K=0 and the two closed forms) moved to
%     offset-mechanism-derivation.tex, which this file does not depend on.
% v1: contained both; the mixture obscured the simple statement.
%
% THIS IS A FRAGMENT, not a standalone document. It uses \mhr from
% jan-augmentation-writeup.tex and is meant to be pulled in with
%     \input{offset-mechanism-equations}
% Preview it with offset-mechanism-preview.tex. Requires amsmath.
%
% Supports SHOWCASE_offset_mechanism_<record> in scripts/gantry/msd-offset/figures/.

\subsection{What the baseline is missing, as two corrections}

The baseline treats the payload as a single rigid mass $m_h$ sitting at $Y$. In reality it is
two things: a rigid part $\mhr$ at $Y$, and an absorber mass $m_a$ on a spring whose rest
position is a distance $L_0$ further out, at $Y+L_0$. The split conserves the total,
$m_h = \mhr + m_a$. Two independent corrections follow, and neither touches the damping or
stiffness matrices, so those are the same in both models throughout.

\paragraph{Starting point.} Of the baseline's three equations, the two that matter here are
%
\begin{align}
  X:\quad & m_t\,\ddot X + M_{12}\,\ddot\Theta + (\text{damping}) = F_X , \label{eq:fix-baseX}\\
  Y:\quad & m_h\,\ddot Y - m_h d\,\ddot\Theta + c_y \dot Y = F_Y , \label{eq:fix-baseY}
\end{align}
%
where $M_{12} = \tfrac{(m_1-m_2)L_b}{2} - m_h Y$ is the coefficient that couples twisting to
sideways motion.

\paragraph{Fix 1: put the payload's mass where it really is.} Wherever the payload's position
enters the mass matrix, substitute the correctly split moments,
%
\begin{equation}
  m_h Y \;\longrightarrow\; \mhr Y + m_a(Y{+}L_0) ,
  \qquad
  m_h Y^2 \;\longrightarrow\; \mhr Y^2 + m_a(Y{+}L_0)^2 .
  \label{eq:fix1-substitution}
\end{equation}
%
Only two coefficients move as a result,
%
\begin{equation}
  M_{12} \;\longrightarrow\; M_{12} - m_a L_0 ,
  \qquad
  M_{22} \;\longrightarrow\; M_{22} + m_a\bigl[2YL_0 + L_0^2\bigr] .
  \label{eq:fix1}
\end{equation}
%
Nothing is added: the model keeps its three coordinates and three equations, and two numbers are
corrected. The size of the correction is a first moment of mass,
$m_a L_0 = 0.101\,\mathrm{kg\,m}$, which divided by $m_h = 10.1\,\mathrm{kg}$ is a
$1\,\mathrm{cm}$ error in the assumed position of the payload's centre of mass.

\paragraph{Fix 2: give the model the wobbling mass.} Add one term to the $Y$ equation and one
new equation for the absorber coordinate $\delta_a$,
%
\begin{align}
  Y:\quad & m_h \ddot Y + \underbrace{m_a \ddot\delta_a}_{\text{new}}
            - m_h d\,\ddot\Theta + c_y \dot Y = F_Y , \label{eq:fix2-y}\\[1mm]
  \delta_a:\quad & m_a\bigl(\ddot Y + \ddot\delta_a - d\,\ddot\Theta\bigr)
            + c_a \dot\delta_a + k_a \delta_a = 0 . \label{eq:fix2-absorber}
\end{align}
%
The new equation is just a mass on a spring: $\ddot Y + \ddot\delta_a - d\,\ddot\Theta$ is the
absorber's acceleration in space, and its inertia is balanced by its own spring $k_a$ and damper
$c_a$. No external force reaches it, hence the zero on the right.

\paragraph{Why each fixes one axis.} The two corrections act on disjoint equations:
%
\begin{center}
\begin{tabular}{@{}lll@{}}
  Fix 1, eq.~\eqref{eq:fix1} & changes rows $X$ and $\Theta$ only
    & so the $Y$ row is untouched and $Y$ is unaffected \\
  Fix 2, eq.~\eqref{eq:fix2-y}--\eqref{eq:fix2-absorber} & changes rows $Y$ and $\delta_a$ only
    & so the $X$ row is untouched and $X$ is unaffected \\
\end{tabular}
\end{center}
%
Neither correction can reach the other's axis. This is why removing one of them cancels the
offset on exactly one output and leaves the other unchanged to six digits, rather than
approximately: the dissociation is structural, not numerical.

Both corrections do perturb the $\Theta$ row, through $M_{22}$ and through
$-m_a d\,\ddot\Theta$ respectively, yet $\Theta$ develops no offset. The reason is that $\Theta$
is the only sprung axis, $K_{22} = k_{b1}{+}k_{b2} > 0$, while $K_{11} = K_{33} = 0$: an error
on $X$ or $Y$ has nothing to restore it and accumulates, whereas on $\Theta$ it decays. That
contrast is what identifies the mechanism, and it is the control row of the figure.

% Preview it with offset-mechanism-preview.tex. Requires amsmath.
%
% Supports SHOWCASE_offset_mechanism_<record> in scripts/gantry/msd-offset/figures/.

\subsection{What the baseline is missing, as two corrections}

The baseline treats the payload as a single rigid mass $m_h$ sitting at $Y$. In reality it is
two things: a rigid part $\mhr$ at $Y$, and an absorber mass $m_a$ on a spring whose rest
position is a distance $L_0$ further out, at $Y+L_0$. The split conserves the total,
$m_h = \mhr + m_a$. Two independent corrections follow, and neither touches the damping or
stiffness matrices, so those are the same in both models throughout.

\paragraph{Starting point.} Of the baseline's three equations, the two that matter here are
%
\begin{align}
  X:\quad & m_t\,\ddot X + M_{12}\,\ddot\Theta + (\text{damping}) = F_X , \label{eq:fix-baseX}\\
  Y:\quad & m_h\,\ddot Y - m_h d\,\ddot\Theta + c_y \dot Y = F_Y , \label{eq:fix-baseY}
\end{align}
%
where $M_{12} = \tfrac{(m_1-m_2)L_b}{2} - m_h Y$ is the coefficient that couples twisting to
sideways motion.

\paragraph{Fix 1: put the payload's mass where it really is.} Wherever the payload's position
enters the mass matrix, substitute the correctly split moments,
%
\begin{equation}
  m_h Y \;\longrightarrow\; \mhr Y + m_a(Y{+}L_0) ,
  \qquad
  m_h Y^2 \;\longrightarrow\; \mhr Y^2 + m_a(Y{+}L_0)^2 .
  \label{eq:fix1-substitution}
\end{equation}
%
Only two coefficients move as a result,
%
\begin{equation}
  M_{12} \;\longrightarrow\; M_{12} - m_a L_0 ,
  \qquad
  M_{22} \;\longrightarrow\; M_{22} + m_a\bigl[2YL_0 + L_0^2\bigr] .
  \label{eq:fix1}
\end{equation}
%
Nothing is added: the model keeps its three coordinates and three equations, and two numbers are
corrected. The size of the correction is a first moment of mass,
$m_a L_0 = 0.101\,\mathrm{kg\,m}$, which divided by $m_h = 10.1\,\mathrm{kg}$ is a
$1\,\mathrm{cm}$ error in the assumed position of the payload's centre of mass.

\paragraph{Fix 2: give the model the wobbling mass.} Add one term to the $Y$ equation and one
new equation for the absorber coordinate $\delta_a$,
%
\begin{align}
  Y:\quad & m_h \ddot Y + \underbrace{m_a \ddot\delta_a}_{\text{new}}
            - m_h d\,\ddot\Theta + c_y \dot Y = F_Y , \label{eq:fix2-y}\\[1mm]
  \delta_a:\quad & m_a\bigl(\ddot Y + \ddot\delta_a - d\,\ddot\Theta\bigr)
            + c_a \dot\delta_a + k_a \delta_a = 0 . \label{eq:fix2-absorber}
\end{align}
%
The new equation is just a mass on a spring: $\ddot Y + \ddot\delta_a - d\,\ddot\Theta$ is the
absorber's acceleration in space, and its inertia is balanced by its own spring $k_a$ and damper
$c_a$. No external force reaches it, hence the zero on the right.

\paragraph{Why each fixes one axis.} The two corrections act on disjoint equations:
%
\begin{center}
\begin{tabular}{@{}lll@{}}
  Fix 1, eq.~\eqref{eq:fix1} & changes rows $X$ and $\Theta$ only
    & so the $Y$ row is untouched and $Y$ is unaffected \\
  Fix 2, eq.~\eqref{eq:fix2-y}--\eqref{eq:fix2-absorber} & changes rows $Y$ and $\delta_a$ only
    & so the $X$ row is untouched and $X$ is unaffected \\
\end{tabular}
\end{center}
%
Neither correction can reach the other's axis. This is why removing one of them cancels the
offset on exactly one output and leaves the other unchanged to six digits, rather than
approximately: the dissociation is structural, not numerical.

Both corrections do perturb the $\Theta$ row, through $M_{22}$ and through
$-m_a d\,\ddot\Theta$ respectively, yet $\Theta$ develops no offset. The reason is that $\Theta$
is the only sprung axis, $K_{22} = k_{b1}{+}k_{b2} > 0$, while $K_{11} = K_{33} = 0$: an error
on $X$ or $Y$ has nothing to restore it and accumulates, whereas on $\Theta$ it decays. That
contrast is what identifies the mechanism, and it is the control row of the figure.

% Preview it with offset-mechanism-preview.tex. Requires amsmath.
%
% Supports SHOWCASE_offset_mechanism_<record> in scripts/gantry/msd-offset/figures/.

\subsection{What the baseline is missing, as two corrections}

The baseline treats the payload as a single rigid mass $m_h$ sitting at $Y$. In reality it is
two things: a rigid part $\mhr$ at $Y$, and an absorber mass $m_a$ on a spring whose rest
position is a distance $L_0$ further out, at $Y+L_0$. The split conserves the total,
$m_h = \mhr + m_a$. Two independent corrections follow, and neither touches the damping or
stiffness matrices, so those are the same in both models throughout.

\paragraph{Starting point.} Of the baseline's three equations, the two that matter here are
%
\begin{align}
  X:\quad & m_t\,\ddot X + M_{12}\,\ddot\Theta + (\text{damping}) = F_X , \label{eq:fix-baseX}\\
  Y:\quad & m_h\,\ddot Y - m_h d\,\ddot\Theta + c_y \dot Y = F_Y , \label{eq:fix-baseY}
\end{align}
%
where $M_{12} = \tfrac{(m_1-m_2)L_b}{2} - m_h Y$ is the coefficient that couples twisting to
sideways motion.

\paragraph{Fix 1: put the payload's mass where it really is.} Wherever the payload's position
enters the mass matrix, substitute the correctly split moments,
%
\begin{equation}
  m_h Y \;\longrightarrow\; \mhr Y + m_a(Y{+}L_0) ,
  \qquad
  m_h Y^2 \;\longrightarrow\; \mhr Y^2 + m_a(Y{+}L_0)^2 .
  \label{eq:fix1-substitution}
\end{equation}
%
Only two coefficients move as a result,
%
\begin{equation}
  M_{12} \;\longrightarrow\; M_{12} - m_a L_0 ,
  \qquad
  M_{22} \;\longrightarrow\; M_{22} + m_a\bigl[2YL_0 + L_0^2\bigr] .
  \label{eq:fix1}
\end{equation}
%
Nothing is added: the model keeps its three coordinates and three equations, and two numbers are
corrected. The size of the correction is a first moment of mass,
$m_a L_0 = 0.101\,\mathrm{kg\,m}$, which divided by $m_h = 10.1\,\mathrm{kg}$ is a
$1\,\mathrm{cm}$ error in the assumed position of the payload's centre of mass.

\paragraph{Fix 2: give the model the wobbling mass.} Add one term to the $Y$ equation and one
new equation for the absorber coordinate $\delta_a$,
%
\begin{align}
  Y:\quad & m_h \ddot Y + \underbrace{m_a \ddot\delta_a}_{\text{new}}
            - m_h d\,\ddot\Theta + c_y \dot Y = F_Y , \label{eq:fix2-y}\\[1mm]
  \delta_a:\quad & m_a\bigl(\ddot Y + \ddot\delta_a - d\,\ddot\Theta\bigr)
            + c_a \dot\delta_a + k_a \delta_a = 0 . \label{eq:fix2-absorber}
\end{align}
%
The new equation is just a mass on a spring: $\ddot Y + \ddot\delta_a - d\,\ddot\Theta$ is the
absorber's acceleration in space, and its inertia is balanced by its own spring $k_a$ and damper
$c_a$. No external force reaches it, hence the zero on the right.

\paragraph{Why each fixes one axis.} The two corrections act on disjoint equations:
%
\begin{center}
\begin{tabular}{@{}lll@{}}
  Fix 1, eq.~\eqref{eq:fix1} & changes rows $X$ and $\Theta$ only
    & so the $Y$ row is untouched and $Y$ is unaffected \\
  Fix 2, eq.~\eqref{eq:fix2-y}--\eqref{eq:fix2-absorber} & changes rows $Y$ and $\delta_a$ only
    & so the $X$ row is untouched and $X$ is unaffected \\
\end{tabular}
\end{center}
%
Neither correction can reach the other's axis. This is why removing one of them cancels the
offset on exactly one output and leaves the other unchanged to six digits, rather than
approximately: the dissociation is structural, not numerical.

Both corrections do perturb the $\Theta$ row, through $M_{22}$ and through
$-m_a d\,\ddot\Theta$ respectively, yet $\Theta$ develops no offset. The reason is that $\Theta$
is the only sprung axis, $K_{22} = k_{b1}{+}k_{b2} > 0$, while $K_{11} = K_{33} = 0$: an error
on $X$ or $Y$ has nothing to restore it and accumulates, whereas on $\Theta$ it decays. That
contrast is what identifies the mechanism, and it is the control row of the figure.

% Preview it with offset-mechanism-preview.tex. Requires amsmath.
%
% Supports SHOWCASE_offset_mechanism_<record> in scripts/gantry/msd-offset/figures/.

\subsection{What the baseline is missing, as two corrections}

The baseline treats the payload as a single rigid mass $m_h$ sitting at $Y$. In reality it is
two things: a rigid part $\mhr$ at $Y$, and an absorber mass $m_a$ on a spring whose rest
position is a distance $L_0$ further out, at $Y+L_0$. The split conserves the total,
$m_h = \mhr + m_a$. Two independent corrections follow, and neither touches the damping or
stiffness matrices, so those are the same in both models throughout.

\paragraph{Starting point.} Of the baseline's three equations, the two that matter here are
%
\begin{align}
  X:\quad & m_t\,\ddot X + M_{12}\,\ddot\Theta + (\text{damping}) = F_X , \label{eq:fix-baseX}\\
  Y:\quad & m_h\,\ddot Y - m_h d\,\ddot\Theta + c_y \dot Y = F_Y , \label{eq:fix-baseY}
\end{align}
%
where $M_{12} = \tfrac{(m_1-m_2)L_b}{2} - m_h Y$ is the coefficient that couples twisting to
sideways motion.

\paragraph{Fix 1: put the payload's mass where it really is.} Wherever the payload's position
enters the mass matrix, substitute the correctly split moments,
%
\begin{equation}
  m_h Y \;\longrightarrow\; \mhr Y + m_a(Y{+}L_0) ,
  \qquad
  m_h Y^2 \;\longrightarrow\; \mhr Y^2 + m_a(Y{+}L_0)^2 .
  \label{eq:fix1-substitution}
\end{equation}
%
Only two coefficients move as a result,
%
\begin{equation}
  M_{12} \;\longrightarrow\; M_{12} - m_a L_0 ,
  \qquad
  M_{22} \;\longrightarrow\; M_{22} + m_a\bigl[2YL_0 + L_0^2\bigr] .
  \label{eq:fix1}
\end{equation}
%
Nothing is added: the model keeps its three coordinates and three equations, and two numbers are
corrected. The size of the correction is a first moment of mass,
$m_a L_0 = 0.101\,\mathrm{kg\,m}$, which divided by $m_h = 10.1\,\mathrm{kg}$ is a
$1\,\mathrm{cm}$ error in the assumed position of the payload's centre of mass.

\paragraph{Fix 2: give the model the wobbling mass.} Add one term to the $Y$ equation and one
new equation for the absorber coordinate $\delta_a$,
%
\begin{align}
  Y:\quad & m_h \ddot Y + \underbrace{m_a \ddot\delta_a}_{\text{new}}
            - m_h d\,\ddot\Theta + c_y \dot Y = F_Y , \label{eq:fix2-y}\\[1mm]
  \delta_a:\quad & m_a\bigl(\ddot Y + \ddot\delta_a - d\,\ddot\Theta\bigr)
            + c_a \dot\delta_a + k_a \delta_a = 0 . \label{eq:fix2-absorber}
\end{align}
%
The new equation is just a mass on a spring: $\ddot Y + \ddot\delta_a - d\,\ddot\Theta$ is the
absorber's acceleration in space, and its inertia is balanced by its own spring $k_a$ and damper
$c_a$. No external force reaches it, hence the zero on the right.

\paragraph{Why each fixes one axis.} The two corrections act on disjoint equations:
%
\begin{center}
\begin{tabular}{@{}lll@{}}
  Fix 1, eq.~\eqref{eq:fix1} & changes rows $X$ and $\Theta$ only
    & so the $Y$ row is untouched and $Y$ is unaffected \\
  Fix 2, eq.~\eqref{eq:fix2-y}--\eqref{eq:fix2-absorber} & changes rows $Y$ and $\delta_a$ only
    & so the $X$ row is untouched and $X$ is unaffected \\
\end{tabular}
\end{center}
%
Neither correction can reach the other's axis. This is why removing one of them cancels the
offset on exactly one output and leaves the other unchanged to six digits, rather than
approximately: the dissociation is structural, not numerical.

Both corrections do perturb the $\Theta$ row, through $M_{22}$ and through
$-m_a d\,\ddot\Theta$ respectively, yet $\Theta$ develops no offset. The reason is that $\Theta$
is the only sprung axis, $K_{22} = k_{b1}{+}k_{b2} > 0$, while $K_{11} = K_{33} = 0$: an error
on $X$ or $Y$ has nothing to restore it and accumulates, whereas on $\Theta$ it decays. That
contrast is what identifies the mechanism, and it is the control row of the figure.
